% Template: templates/main.tex (the user-supplied template for this run). Structure preserved; content tailored.
% Target: PwC Acceleration Center India — Business AI Engineer - Intern - 2027.
\documentclass[10pt, letterpaper]{article}

% Packages:
\usepackage[
    ignoreheadfoot, % set margins without considering header and footer
    top=2 cm, % seperation between body and page edge from the top
    bottom=2 cm, % seperation between body and page edge from the bottom
    left=2 cm, % seperation between body and page edge from the left
    right=2 cm, % seperation between body and page edge from the right
    footskip=1.0 cm, % seperation between body and footer
    % showframe % for debugging
]{geometry} % for adjusting page geometry
\usepackage{titlesec} % for customizing section titles
\usepackage{tabularx} % for making tables with fixed width columns
\usepackage{array} % tabularx requires this
\usepackage[dvipsnames]{xcolor} % for coloring text
\definecolor{primaryColor}{RGB}{0, 0, 0} % define primary color
\usepackage{enumitem} % for customizing lists
\usepackage{fontawesome5} % for using icons
\usepackage{amsmath} % for math
\usepackage[
    pdftitle={Soham Datta — Resume},
    pdfauthor={Soham Datta},
    pdfcreator={LaTeX},
    colorlinks=true,
    urlcolor=primaryColor
]{hyperref} % for links, metadata and bookmarks
\usepackage[pscoord]{eso-pic} % for floating text on the page
\usepackage{calc} % for calculating lengths
\usepackage{bookmark} % for bookmarks
\usepackage{lastpage} % for getting the total number of pages
\usepackage{changepage} % for one column entries (adjustwidth environment)
\usepackage{paracol} % for two and three column entries
\usepackage{ifthen} % for conditional statements
\usepackage{needspace} % for avoiding page brake right after the section title
\usepackage{iftex} % check if engine is pdflatex, xetex or luatex

% Ensure that generate pdf is machine readable/ATS parsable:
\ifPDFTeX
    \input{glyphtounicode}
    \pdfgentounicode=1
    \usepackage[T1]{fontenc}
    \usepackage[utf8]{inputenc}
    \usepackage{lmodern}
\fi

\usepackage{charter}

% Some settings:
\raggedright
\AtBeginEnvironment{adjustwidth}{\partopsep0pt} % remove space before adjustwidth environment
\pagestyle{empty} % no header or footer
\setcounter{secnumdepth}{0} % no section numbering
\setlength{\parindent}{0pt} % no indentation
\setlength{\topskip}{0pt} % no top skip
\setlength{\columnsep}{0.15cm} % set column seperation
\pagenumbering{gobble} % no page numbering

\titleformat{\section}{\needspace{4\baselineskip}\bfseries\large}{}{0pt}{}[\vspace{1pt}\titlerule]

\titlespacing{\section}{
    % left space:
    -1pt
}{
    % top space:
    0.16 cm
}{
    % bottom space:
    0.12 cm
} % section title spacing

\renewcommand\labelitemi{$\vcenter{\hbox{\small$\bullet$}}$} % custom bullet points
\newenvironment{highlights}{
    \begin{itemize}[
        topsep=0.10 cm,
        parsep=0.10 cm,
        partopsep=0pt,
        itemsep=0pt,
        leftmargin=0 cm + 10pt
    ]
}{
    \end{itemize}
} % new environment for highlights


\newenvironment{highlightsforbulletentries}{
    \begin{itemize}[
        topsep=0.10 cm,
        parsep=0.10 cm,
        partopsep=0pt,
        itemsep=0pt,
        leftmargin=10pt
    ]
}{
    \end{itemize}
} % new environment for highlights for bullet entries

\newenvironment{onecolentry}{
    \begin{adjustwidth}{
        0 cm + 0.00001 cm
    }{
        0 cm + 0.00001 cm
    }
}{
    \end{adjustwidth}
} % new environment for one column entries

\newenvironment{twocolentry}[2][]{
    \onecolentry
    \def\secondColumn{#2}
    \setcolumnwidth{\fill, 4.5 cm}
    \begin{paracol}{2}
}{
    \switchcolumn \raggedleft \secondColumn
    \end{paracol}
    \endonecolentry
} % new environment for two column entries

\newenvironment{threecolentry}[3][]{
    \onecolentry
    \def\thirdColumn{#3}
    \setcolumnwidth{, \fill, 4.5 cm}
    \begin{paracol}{3}
    {\raggedright #2} \switchcolumn
}{
    \switchcolumn \raggedleft \thirdColumn
    \end{paracol}
    \endonecolentry
} % new environment for three column entries

\newenvironment{header}{
    \setlength{\topsep}{0pt}\par\kern\topsep\centering\linespread{1.5}
}{
    \par\kern\topsep
} % new environment for the header

% save the original href command in a new command:
\let\hrefWithoutArrow\href

% new command for external links:


\begin{document}
    \newcommand{\AND}{\unskip
        \cleaders\copy\ANDbox\hskip\wd\ANDbox
        \ignorespaces
    }
    \newsavebox\ANDbox
    \sbox\ANDbox{$|$}

    \begin{header}
        \fontsize{25 pt}{25 pt}\selectfont Soham Datta

        \vspace{5 pt}

        \normalsize
        \mbox{Pune, India}%
        \kern 5.0 pt%
        \AND%
        \kern 5.0 pt%
        \mbox{\hrefWithoutArrow{mailto:thesohamdatta@gmail.com}{thesohamdatta@gmail.com}}%
        \kern 5.0 pt%
        \AND%
        \kern 5.0 pt%
        \mbox{+91 9420984066}%
        \kern 5.0 pt%
        \AND%
        \kern 5.0 pt%
        \mbox{\hrefWithoutArrow{https://www.linkedin.com/in/thesohamdatta/}{linkedin.com/in/thesohamdatta}}%
        \kern 5.0 pt%
        \AND%
        \kern 5.0 pt%
        \mbox{\hrefWithoutArrow{https://github.com/thesohamdatta}{github.com/thesohamdatta}}%
        \kern 5.0 pt%
        \AND%
        \kern 5.0 pt%
        \mbox{\hrefWithoutArrow{https://sohamdatta.framer.ai}{sohamdatta.framer.ai}}%
    \end{header}

    \vspace{5 pt - 0.3 cm}

    \begin{onecolentry}
        {\small\bfseries AI Engineer Intern | AI Fluency, Data Analysis, Digital Tooling}
    \end{onecolentry}

    \vspace{0.05 cm}

    \begin{onecolentry}
        Final-year AI \& ML undergraduate (B.E., expected Jun 2026) building LLM applications and agentic tooling that turn unstructured information into structured, reviewable output, with hands-on computer-vision and research-data work.
    \end{onecolentry}

    \section{Experience}

        \begin{twocolentry}{
            Jun 2026 -- Present
        }
            \textbf{Founder}, Aura\end{twocolentry}

        \vspace{0.06 cm}
        \begin{onecolentry}
            \begin{highlights}
                \item Built an ESP32-S3 wearable AI system end to end: C/C++ firmware for 16 kHz I2S DMA audio capture, Opus compression, FreeRTOS event handling, and OTA updates, plus FastAPI services for streaming audio, memory/context extraction, and retrieval over a Supabase/pgvector vector store.
                \item Adapted and extended Omi open-source components for knowledge-graph linking and speaker identification, and ran a Kubernetes-hosted Deepgram ASR integration.
            \end{highlights}
        \end{onecolentry}

        \vspace{0.12 cm}

        \begin{twocolentry}{
            Oct 2024 -- Dec 2024
        }
            \textbf{AI/ML Intern}, AICTE\end{twocolentry}

        \vspace{0.06 cm}
        \begin{onecolentry}
            \begin{highlights}
                \item Built and evaluated on-device computer-vision workflows in TensorFlow, TensorFlow Lite, and Google ML Kit for custom object detection (Ef\mbox{f}icientDet-Lite) and image classification on Android.
            \end{highlights}
        \end{onecolentry}

        \vspace{0.12 cm}

        \begin{twocolentry}{
            Jul 2024 -- Oct 2024
        }
            \textbf{Associate}, Reliance\end{twocolentry}

        \vspace{0.06 cm}
        \begin{onecolentry}
            \begin{highlights}
                \item Handled customer-service and store operations across POS, memberships, vouchers, and exchanges, keeping daily documentation and reporting accurate and coordinating with managers and store teams.
            \end{highlights}
        \end{onecolentry}

        \vspace{0.12 cm}

        \begin{twocolentry}{
            Sep 2023 -- Mar 2024
        }
            \textbf{Research / Library Assistant}, PES Modern College of Engineering\end{twocolentry}

        \vspace{0.06 cm}
        \begin{onecolentry}
            \begin{highlights}
                \item Supported literature discovery, research documentation, and digitization, organising 3,000+ pages of source material into a searchable archive for students and staff.
            \end{highlights}
        \end{onecolentry}

    \section{Projects}

        \begin{twocolentry}{
            {\footnotesize\href{https://github.com/thesohamdatta/Mia}{github.com/thesohamdatta/Mia}}
        }
            \textbf{Mia: AI Engineering CLI} -- TypeScript, Bun, LLM adapters, Vitest\end{twocolentry}

        \vspace{0.06 cm}
        \begin{onecolentry}
            \begin{highlights}
                \item Built a TypeScript CLI that runs a five-stage AI engineering workflow (grill, spec, plan, review, ship) behind multi-provider LLM adapters for Claude, Codex, and local models.
                \item Kept state in an immutable JSONL event log and gated every stage with automated lint, type-check, and Vitest checks, using stateless, event-driven execution (ADR-0001).
            \end{highlights}
        \end{onecolentry}

        \vspace{0.12 cm}

        \begin{twocolentry}{
            {\footnotesize\href{https://github.com/thesohamdatta/LLM-Council}{github.com/thesohamdatta/LLM-Council}}
        }
            \textbf{LLM-Council: Multi-Model Evaluation Service} -- Python, FastAPI, React\end{twocolentry}

        \vspace{0.06 cm}
        \begin{onecolentry}
            \begin{highlights}
                \item Built a Python FastAPI service and React client that query several LLM providers, then compare and synthesize their answers through an Analyst, Skeptic, Synthesizer pipeline.
            \end{highlights}
        \end{onecolentry}

        \vspace{0.12 cm}

        \begin{twocolentry}{
            {\footnotesize\href{https://github.com/BasedHardware/omi}{github.com/BasedHardware/omi}}
        }
            \textbf{Omi (Open Source): Wearable AI Platform} -- Contributor\end{twocolentry}

        \vspace{0.06 cm}
        \begin{onecolentry}
            \begin{highlights}
                \item Authored offline speaker clustering with 52 unit tests (approach adopted downstream), and diagnosed a Windows Electron/WebGL GPU issue in a filed root-cause report.
            \end{highlights}
        \end{onecolentry}

    \section{Education}

        \begin{twocolentry}{
            2022 -- Jun 2026
        }
            \textbf{Savitribai Phule Pune University (SPPU)}, B.E. in Artificial Intelligence \& Machine Learning -- Pune, India\end{twocolentry}

    \section{Technical Skills}

        \begin{onecolentry}
            \textbf{Languages:} Python, TypeScript/JavaScript, Kotlin, C/C++
        \end{onecolentry}

        \vspace{0.12 cm}

        \begin{onecolentry}
            \textbf{AI \& Data:} LLM applications, AI agents, RAG, tool calling, multi-agent systems, computer vision, TensorFlow, TensorFlow Lite, PyTorch, Google ML Kit
        \end{onecolentry}

        \vspace{0.12 cm}

        \begin{onecolentry}
            \textbf{Digital Tools \& Engineering:} FastAPI, React, REST APIs, Kubernetes, Supabase/pgvector, Git/GitHub, Vitest, documentation and reporting
        \end{onecolentry}
\end{document}
